\section{Appendix A: Observable estimation and supplied-graph comparison}

The observable upper guarantee rests on a separation of treatment-assignment variation from edge-audit variation. The audit-score moment identity makes that separation exact, while the supplied-graph comparison identifies a component of minimax risk that persists throughout the retention range. The correspondence with the published additive model then fixes the degree convention and observation channel needed to interpret these comparisons.

\subsection{Audit-score moments and observable estimation}

For a fixed schedule, the audit score in \cref{obj:def:audit-score} weights recorded spillover coordinates by their inverse inclusion probability. Its assignment-variance benchmark is the oracle score in \cref{obj:def:oracle-score}, which uses the true graph. The following moment statement connects these scores to the population causal target.

\begin{lemmav}{lem:score-audit}[Audit-score moments and oracle bound]
Let \(V\) be a finite population with \(n=|V|\), and let \(\theta=(G,a,t,b)\) be a fixed schedule. Write \(Y_i(z)\) for unit \(i\)'s additive potential outcome under binary assignment \(z\). Suppose:
\begin{itemize}
\item \textbf{(Population and degree.)} \(n\ge4\), and \(d\) is an integer satisfying \(1\le d\le n-1\).
\item \textbf{(Schedule class.)} \(\theta\in\mathcal M_{n,d}\) as in \cref{obj:def:model}.
\item \textbf{(Assignment.)} The treatment-assignment vector \(Z\) satisfies \cref{obj:ass:assignment}.
\item \textbf{(Audit.)} The off-diagonal audit-mark array \(W\) satisfies \cref{obj:ass:audit} with retention probability \(q>0\).
\item \textbf{(Design independence.)} \(Z\) and \(W\) satisfy \cref{obj:ass:design-independence}.
\end{itemize}
Let \(P_{\theta,q}\) be the induced experiment law for this fixed schedule, with the observed record retaining precisely the audited arrows of \(G\), the assignment, and the realized outcomes. Let \(T_q\) and \(T_G\) be the scores in \cref{obj:def:audit-score,obj:def:oracle-score}, respectively, and define the all-treated-versus-all-control population contrast by
\[
\tau(\theta)
=\frac1n\sum_{i\in V}
\bigl(Y_i(\mathbf1)-Y_i(\mathbf0)\bigr),
\]
where \(\mathbf1\) and \(\mathbf0\) are the all-treated and all-control assignments. Then
\[
\mathbb E_{P_{\theta,q}}[T_q]=\tau(\theta),
\]
\[
\operatorname{Var}_{P_{\theta,q}}(T_q)
=
\operatorname{Var}_{Z}(T_G)
+\frac{4(1-q)}{n^2q}
\sum_{i\in V}
\bigl|\{j\in V:(j,i)\in G\}\bigr|
\,\mathbb E_Z\!\left[Y_i(Z)^2\right],
\]
and
\[
\operatorname{Var}_{Z}(T_G)\le\frac{5(d+1)^2}{n},
\]
where expectations and variances indexed by \(Z\) use the independent half-Bernoulli assignment law.
\end{lemmav}

The unbiasedness statement concerns expectation over both parts of the experimental design for each fixed schedule. In the variance identity, the oracle term measures assignment variation, and the additional term measures the variation introduced by independently retaining true arrows. Each recipient's contribution to this additional term is weighted by its true in-degree and its assignment-averaged squared outcome. Thus the identity distinguishes the variance due to recording edges from the variance present when the graph is supplied.

The coefficient-mass restriction in \cref{obj:ass:coefficient-mass} controls both components of the upper-risk calculation. It bounds every potential outcome in absolute value by one and places the population contrast in \([-1,1]\). Together with the in-degree bound, it bounds the audit contribution by \(4d(1-q)/(nq)\). The common in-degree and out-degree restrictions enter the oracle bound in \cref{obj:lem:score-audit}, which is uniform over the schedule class. These two contributions are the terms in the upper envelope of \cref{obj:def:upper-envelope}.

The observable guarantee in \cref{obj:thm:observable-upper} uses the total rule of \cref{obj:def:upper-rule}. At positive retention and an upper envelope below one, this rule projects the audit score onto \([-1,1]\); projection preserves its squared-error upper bound because the causal target lies in that interval. For the remaining parameter values, the zero branch has squared loss at most one. Measurability follows from the score's finite sums and products of observed graph, assignment, and outcome coordinates, together with the continuous clipping map. The branch choice depends on the known design parameters. Consequently, the rule belongs to the estimator class in \cref{obj:def:risk}, including at zero retention.

\subsection{The supplied-graph lower comparison}

The oracle variance bound gives an attainable precision benchmark when the true graph is supplied. A complementary minimax statement establishes the degree-dependent lower scale for both observation experiments. The complete-record testing arguments in \cref{obj:lem:full-record-two-prior-risk} provide the risk-reduction framework for this comparison.

\begin{lemmav}{lem:supplied-block-record-lower}[Complete-record minimax lower bounds]
Let \(n\) be the population size, \(d\) the common off-diagonal in-degree and out-degree bound, and \(q\) the known edge-retention probability. Suppose that:
\begin{itemize}
\item \textbf{(Parameter range.)} \(n,d\) are integers with \(n\ge4\) and \(1\le d\le n-1\), and \(q\in[0,1]\).
\item \textbf{(Assignment law.)} The joint law of the treatment assignment and audit marks satisfies \cref{obj:ass:assignment}.
\item \textbf{(Audit law.)} The same joint law satisfies \cref{obj:ass:audit} at retention probability \(q\).
\item \textbf{(Design independence.)} The same joint law satisfies \cref{obj:ass:design-independence}.
\end{itemize}
Then the complete-record minimax squared risk \(R_{n,d}(q)\) of \cref{obj:def:risk} and the supplied-graph minimax squared risk \(R^G_{n,d}(q)\) of \cref{obj:def:supplied-risk} satisfy
\[
R_{n,d}(q)\ge
\frac{1}{160000\pi^2}
\min\left\{1,\frac{(d+1)^2}{n}\right\},
\qquad
R^G_{n,d}(q)\ge
\frac{1}{160000\pi^2}
\min\left\{1,\frac{(d+1)^2}{n}\right\}.
\]
These bounds apply to every measurable randomized estimator in the respective definitions, with the true off-diagonal graph supplied in the second experiment.
\end{lemmav}

Both lower bounds are uniform in the retention probability and cover arbitrary measurable use of the available record and estimator randomization. In particular, the supplied-graph bound establishes the degree-dependent scale even when every true arrow is available to the estimator. Combined with the oracle bound in \cref{obj:lem:score-audit} and clipping to the target interval, it characterizes supplied-graph minimax precision up to universal constants. The audit-score variance identity then quantifies the additional variance of the observable score relative to that benchmark.

\subsection{Published-model correspondence}

A comparison with the published SNIPE model requires matching its graph and coefficient conventions. In the interaction-order-one specification of \citet[Sections 3.1--3.2, Assumptions 1--2, equations (3.2)--(3.3)]{CortezRodriguezEichhornYu2023SNIPE}, outcomes are polynomials with a constant term and singleton-coordinate terms indexed by each recipient's in-neighborhood. Those neighborhoods may contain the recipient itself, and the coefficient-mass envelope includes the constant coefficient. The following correspondence therefore adjoins each unit's own coordinate to the off-diagonal graph used here.

\begin{propositionv}{lem:published-model}[Published-model correspondence and risk transfer]*
Let \(V\) be the fixed finite population. Use the published SNIPE model conventions of \citep[Sections 3.1--3.2, Assumptions 1--2, equations (3.2)--(3.3)]{CortezRodriguezEichhornYu2023SNIPE}. For every additive schedule \(\theta=(G,a,t,b)\), define its augmented graph \(\widetilde G\) and published polynomial coefficients by
\[
\widetilde G=G\cup\{(i,i):i\in V\},\qquad
c^{\mathrm{pub}}_{i,\varnothing}=a_i,\qquad
c^{\mathrm{pub}}_{i,\{i\}}=t_i,\qquad
c^{\mathrm{pub}}_{i,\{j\}}=b_{ij}\quad ((j,i)\in G),
\]
with all remaining polynomial coefficients zero.

For every nonnegative integer \(d\), if \(\theta\in\mathcal M_{n,d}\), the schedule class in \cref{obj:def:model}, its augmentation has every diagonal arrow and belongs to the published specialization with interaction order \(\beta=1\), coefficient-mass envelope
\[
Y_{\max}
=\max_{i\in V}
\left(
|c^{\mathrm{pub}}_{i,\varnothing}|
+\sum_{j:(j,i)\in\widetilde G}
|c^{\mathrm{pub}}_{i,\{j\}}|
\right)
\le 1,
\]
and total in-degree and out-degree at most \(d+1\). Conversely, every published schedule satisfying these order, envelope, and degree restrictions and containing every diagonal arrow yields a member of \(\mathcal M_{n,d}\) by deleting the diagonal arrows and reading its empty-set coefficient as \(a_i\), its own singleton coefficient as \(t_i\), and its singleton coefficient for source \(j\) as \(b_{ij}\). Augmenting this stripped schedule recovers the original graph and every polynomial coefficient indexed by a subset of a recipient's in-neighborhood. For every additive schedule, stripping its augmentation recovers its graph, baseline vector, own-treatment vector, and spillover coefficients on all graph arrows.

For every additive schedule, every unit \(i\), and every binary assignment \(z\), augmentation preserves \(Y_i(z)\), the additive potential outcome, and preserves \(\tau(\theta)\), the population all-treated-versus-all-control contrast. For every assignment and audit-mark array, adjoining the known diagonal arrows to \(H\), the retained graph, gives the published record through the deterministic transformation
\[
O=(H,Z,Y(Z))
\longmapsto
\bigl(H\cup\{(i,i):i\in V\},\,Z,\,Y(Z)\bigr),
\]
where \(O\) is the complete observed record, \(Z\) is the assignment vector, and \(Y(Z)\) is the realized-outcome vector.

For every joint assignment-and-audit law satisfying \cref{obj:ass:assignment,obj:ass:audit,obj:ass:design-independence}, with \(q\) the edge-retention probability, and every nonnegative integer \(d\), consider any class of published schedules containing the augmentation of every member of \(\mathcal M_{n,d}\). Under the same assignment law and off-diagonal audit channel, \(R_{n,d}(q)\), the minimax squared risk defined in \cref{obj:def:risk}, is at most the minimax squared risk over that published class.
\end{propositionv}
% CAUSALSMITH-CITED-SCOPE-BEGIN lem:published-model
\verificationfootnotetext{\textbf{Formalization scope.} This result uses the cited conclusion \citep{CortezRodriguezEichhornYu2023SNIPE} (Sections 3.1--3.2, Assumptions 1--2 and equations (3.2)--(3.3), arXiv version 4.). The cited source proof is not formalized here: Lean verifies its use and all remaining steps. If that cited conclusion is false, the portions of this result that depend on it are not certified.}
% CAUSALSMITH-CITED-SCOPE-END lem:published-model

The augmented graph \(\widetilde G\) in \cref{obj:lem:published-model} has one additional incoming and outgoing arrow per unit. Accordingly, the published total-degree index is \(d+1\) when the off-diagonal degree index is \(d\). The own-treatment coefficient occupies the recipient's singleton coordinate, and the published coefficient-mass envelope agrees with the row restriction in \cref{obj:ass:coefficient-mass}. The cited model conventions permit zero coefficients on graph arrows, so the correspondence retains the declared graph together with its coefficients.

The risk-transfer clause uses both schedule inclusion and a common observation channel. Under the same assignment law and off-diagonal audit mechanism, adjoining the known diagonal arrows transforms the observed record while preserving realized outcomes and the causal target. The resulting risk ordering transfers a lower bound from the additive class to any published class containing all of its augmentations. Comparisons with a supplied true graph use the separate experiment in \cref{obj:def:supplied-risk}; the published-model correspondence specifies the channel under which its own class-inclusion comparison applies.

\subsection{Supported boundary cases}

The observable upper guarantee and supplied-graph lower comparison fix the full-retention scale. The complete-record testing bounds in \cref{obj:thm:block-testing-scale,obj:lem:full-record-two-prior-risk} supply the complementary retention-dependent lower scale. The following statement consolidates these ingredients at zero retention, full retention, and degree one.

\begin{theoremv}{thm:supported-boundaries}[Minimax bounds at supported boundaries]
There exist universal constants \(c,C_0>0\) such that the following holds for every population size \(n\), degree bound \(d\), known retention probability \(q\), and joint law of the treatment-assignment vector \(Z\) and the off-diagonal audit-mark array \(W\) satisfying:
\begin{itemize}
\item \textbf{(Population.)} \(n\) is an integer with \(n\ge4\).
\item \textbf{(Degree.)} \(d\) is an integer with \(1\le d\le n-1\).
\item \textbf{(Retention.)} \(q\in[0,1]\).
\item \textbf{(Assignment.)} The joint law satisfies \cref{obj:ass:assignment}.
\item \textbf{(Audit.)} The joint law satisfies \cref{obj:ass:audit} at retention probability \(q\).
\item \textbf{(Independence.)} The joint law satisfies \cref{obj:ass:design-independence}.
\end{itemize}
For each such law, the minimax squared risk \(R_{n,d}(q)\) of \cref{obj:def:risk}, with the infimum taken over all measurable randomized estimators of the complete observed record \(O\), satisfies
\[
R_{n,d}(q)\ge
c\max\left\{
\min\left\{1,\frac{(d+1)^2}{n}\right\},
\frac{1}{1+nq}
\right\}.
\]
The supported boundary cases satisfy
\[
\begin{aligned}
c&\le R_{n,d}(0)\le1,\\
c\min\left\{1,\frac{(d+1)^2}{n}\right\}
&\le R_{n,d}(1)
\le C_0\min\left\{1,\frac{(d+1)^2}{n}\right\},\\
\frac{c}{1+nq}
&\le R_{n,1}(q)\le\frac{C_0}{1+nq}.
\end{aligned}
\]
\end{theoremv}

The two terms in the lower bound express distinct sources of difficulty: the degree-dependent precision scale already present with a supplied graph, and the retention-dependent scale associated with partial observation. At zero retention, the complete-record minimax risk remains bounded below by a universal positive constant. At full retention, the audit contribution in \cref{obj:lem:score-audit} vanishes, and the upper and lower bounds agree up to universal constants at the degree-dependent scale. For degree one, the matching reciprocal scale \(1/(1+nq)\) describes the entire retention range, linking constant risk at sparse collection to order \(1/n\) risk at full retention.
